In vector notation the centroids are $\dfrac{P+A+B}{3},\dfrac{P+A+C}{3},\dfrac{P+B+C}{3}$. Their differences are one third of corresponding side vectors of $ABC$, so the centroid triangle is similar with length factor $\dfrac{1}{3}$ and area factor $\dfrac{1}{9}$. The original right triangle has area $\dfrac{3\cdot4}{2}=6$, giving $\boxed{\dfrac{2}{3}}$. The argument works for every point $P$.
\begin{asy}
import graph;import geometry;usepackage("amsmath");size(7cm);
pair A=(0,3),B=(0,0),C=(4,0),P=(1,1),G1=(P+A+B)/3,G2=(P+A+C)/3,G3=(P+B+C)/3;draw(A--B--C--cycle);draw(P--A^^P--B^^P--C,dashed);draw(G1--G2--G3--cycle,blue);label("$A$",A,(-1,1));label("$B$",B,(-1,-1));label("$C$",C,(1,-1));label("$P$",P,(-1,0));label("$G_1$",G1,(-1,0));label("$G_2$",G2,(1,1));label("$G_3$",G3,(0,-1));
\end{asy}

